Pour des réels $a, b, c, d > 0$, on pose
\[
S(a, b, c, d) = \frac{a}{a+b+d} + \frac{b}{a+b+c} + \frac{c}{b+c+d} + \frac{d}{a+c+d}.
\]
On cherche l'ensemble de toutes les valeurs prises par $S$ (OIM 1974).
\begin{enumerate}
  \item En comparant chaque dénominateur à $a + b + c + d$, montrer que $S > 1$.
  \item Montrer que $\dfrac{a}{a+b+d} + \dfrac{b}{a+b+c} < 1$ (comparer
  chaque dénominateur à $a + b$), puis que $S < 2$.
  \item Calculer $S(t, t, 1, 1)$ et $S(1, t, 1, t)$ pour $t > 0$, puis leurs
  limites quand $t \to 0^+$.
  \item Soit $\sigma \in \left]1, 2\right[$. Justifier l'existence de deux
  quadruplets $P$ et $Q$ tels que $S(P) < \sigma < S(Q)$. En étudiant
  $g(\lambda) = S\bigl((1-\lambda)P + \lambda Q\bigr)$ pour
  $\lambda \in [0, 1]$, montrer que $\sigma$ est une valeur prise par $S$.
  \item Conclure.
\end{enumerate}
